\documentclass[11pt]{article}

\usepackage[a4paper,margin=1in]{geometry}
\usepackage[T1]{fontenc}
\usepackage{lmodern}
\usepackage{microtype}
\usepackage{parskip}
\usepackage{graphicx}
\usepackage{amsmath,amssymb}
\usepackage{booktabs}
\usepackage{array}
\usepackage{caption}
\usepackage{float}
\usepackage{enumitem}
\usepackage{xcolor}
\usepackage{titlesec}
\usepackage{fancyhdr}
\usepackage{listings}
\usepackage[most]{tcolorbox}

\definecolor{primary}{HTML}{1B3A5C}
\definecolor{accent}{HTML}{2E86AB}
\definecolor{accent2}{HTML}{C9622A}
\definecolor{panelbg}{HTML}{F4F6F8}

\renewcommand{\arraystretch}{1.2}

\titleformat{\section}
  {\normalfont\Large\bfseries\color{primary}}
  {\thesection}{0.6em}{}[{\color{accent}\titlerule[1.1pt]}]
\titleformat{\subsection}
  {\normalfont\large\bfseries\color{primary}}
  {\thesubsection}{0.5em}{}
\titleformat{\subsubsection}
  {\normalfont\bfseries\color{accent}}
  {\thesubsubsection}{0.5em}{}
\titlespacing*{\section}{0pt}{1.6em}{0.8em}
\titlespacing*{\subsection}{0pt}{1.2em}{0.5em}

\pagestyle{fancy}
\fancyhf{}
\renewcommand{\headrulewidth}{0.4pt}
\renewcommand{\footrulewidth}{0pt}
\fancyhead[L]{\small\color{primary}\bfseries CS683 PA-1 - Task 1}
\fancyhead[R]{\small\color{primary}\nouppercase{\rightmark}}
\fancyfoot[C]{\small\color{primary}\thepage}

\lstdefinestyle{cmdstyle}{
    backgroundcolor=\color{panelbg},
    basicstyle=\ttfamily\small,
    frame=single,
    rulecolor=\color{accent},
    framerule=0.6pt,
    breaklines=true,
    columns=fullflexible,
    xleftmargin=4pt, xrightmargin=4pt,
    aboveskip=8pt, belowskip=8pt,
}
\lstset{style=cmdstyle}

\newtcolorbox{insightbox}{
    enhanced, breakable, colback=accent!6, colframe=accent,
    boxrule=0.8pt, arc=3pt, left=8pt, right=8pt, top=6pt, bottom=6pt,
    fonttitle=\bfseries, coltitle=white, colbacktitle=accent,
    title={Key Insight}, before skip=8pt, after skip=10pt,
}
\newtcolorbox{todobox}[1][Fill in]{
    enhanced, breakable, colback=accent2!4, colframe=accent2!70!black,
    boxrule=0.6pt, arc=2pt, left=8pt, right=8pt, top=5pt, bottom=5pt,
    fonttitle=\bfseries\small, coltitle=accent2!70!black,
    title={TODO --- #1}, before skip=6pt, after skip=10pt,
}

\usepackage[colorlinks=true,linkcolor=accent,citecolor=accent,urlcolor=accent]{hyperref}

\graphicspath{{../../plots/}}

\begin{document}

\begin{titlepage}
    \centering
    \vspace*{2.2cm}
    {\color{accent}\rule{0.5\textwidth}{1.2pt}}\\[0.6cm]
    {\Huge\bfseries\color{primary} CS683 PA-1}\\[0.4cm]
    {\LARGE\bfseries\color{primary} Task 1 Report}\\[0.3cm]
    {\Large 2D Convolution}\\[0.6cm]
    {\color{accent}\rule{0.5\textwidth}{1.2pt}}\\[1.4cm]

    {\large Loop Reordering \(\cdot\) Loop Unrolling \(\cdot\) Cache Tiling \(\cdot\) SIMD (AVX2)}\\[2cm]

    \vfill
    {\large \today}
\end{titlepage}

\tableofcontents
\newpage

\section*{Setup}
\addcontentsline{toc}{section}{Setup}

All experiments were run on an Intel 13th Gen CPU, which has both
performance (P) cores and efficiency (E) cores. The P-cores have a 48\,KB
L1-D cache, and the E-cores have a 32\,KB L1-D cache. All the measurements
in this report were run on CPU 0, which is a P-core.

\section{Task 1A --- Loop Reordering \& Unrolling}

\subsection{Loop Reordering}
\begin{todobox}[Motivation]
    What motivated the change; what was reordered and why.
\end{todobox}

\subsubsection*{Considerations}
\begin{todobox}[Considerations]
    Unit-stride access pattern; why $K \times K$ passes touch the whole image;
    correctness constraints.
\end{todobox}

\subsubsection*{Results}
\begin{figure}[H]
    \centering
    \caption{Speedup of \texttt{conv\_reorder} over \texttt{conv\_naive} vs.\ image size.}
    \label{fig:reorder-speedup-size}
\end{figure}
\begin{figure}[H]
    \centering
    \caption{Speedup of \texttt{conv\_reorder} over \texttt{conv\_naive} vs.\ kernel size $K$.}
    \label{fig:reorder-speedup-kernel}
\end{figure}

\subsubsection*{Discussion}
\begin{todobox}[Discussion]
    Metric used to quantify effectiveness and why; explain the ``can be as
    slow as or slower than naive at large sizes'' behavior seen at $K=3$.
\end{todobox}

\subsection{Loop Unrolling}
\begin{todobox}[Motivation]
    What was unrolled (the reordered inner loop over output columns) and why
    (expose ILP, hide FP latency).
\end{todobox}

\subsubsection*{Considerations}
\begin{todobox}[Considerations]
    Unroll factor chosen; no remainder tail needed since \texttt{W \% 8 == 0}.
\end{todobox}

\subsubsection*{Results}
\begin{figure}[H]
    \centering
    \caption{Speedup of \texttt{conv\_unroll} over \texttt{conv\_naive} vs.\ image size.}
\end{figure}
\begin{figure}[H]
    \centering
    \caption{Speedup of \texttt{conv\_unroll} over \texttt{conv\_naive} vs.\ kernel size $K$.}
\end{figure}

\subsubsection*{Discussion}
\begin{todobox}[Discussion]
    Interpret the unrolling results in light of Section~2's memory-bound
    behavior.
\end{todobox}

\section{Task 1B --- Tiling}

\subsection{Baseline Profiling}

The CPU we used to run is an Intel hybrid P-core/E-core design. The L1d cache
size, as seen by doing \texttt{cat /sys/devices/system/cpu/cpu0/cache/index0/size},
came out to be 48\,KB for the P-core, and 32\,KB for the E-core. All
measurements below were taken while running on a P-core (CPU 0) with
\texttt{taskset}, so the L1d size for this baseline is \textbf{48\,KB}.

\begin{lstlisting}[language=bash]
taskset -c 0 perf stat -r 5 \
    -e instructions,L1-dcache-loads,L1-dcache-load-misses \
    ./bin/conv naive
\end{lstlisting}

\begin{table}[H]
    \centering
    \begin{tabular}{lr}
        \toprule
        \textbf{Metric} & \textbf{Value (5-run mean)} \\
        \midrule
        Instructions      & 4{,}371{,}265{,}012 \\
        L1d loads         & 892{,}590{,}815 \\
        L1d load misses   & 571{,}072 \\
        Time elapsed (s)  & 0.222255 $\pm$ 0.002410 \\
        \bottomrule
    \end{tabular}
    \caption{\texttt{perf stat} results for \texttt{conv\_naive} on a
    2048$\times$2048, $K=3$ workload, pinned to P-core CPU 0.}
    \label{tab:naive-baseline-perf}
\end{table}

\[
    \text{MPKI} = \frac{\text{L1-D load misses}}{\text{instructions}} \times 1000
    = \frac{571{,}072}{4{,}371{,}265{,}012} \times 1000 \approx 0.131
\]

\begin{insightbox}
    An MPKI of $\approx 0.13$ is pretty low, and it's consistent with what the naive kernel is actually doing. For every
    output pixel, it just needs a tiny $3\times3$ neighborhood of the input
    and the same $3\times3$ kernel it already used one pixel ago. All of
    that easily fits into the cache. So even though we loop through the entire 16\,MB image over the course of the run, at
    any given moment it's only really touching a small, reused patch of
    memory. Hence, we don't get as many misses in this case.
\end{insightbox}

\subsection{Tiled Implementation}

\texttt{conv\_tile} blocks the output into square $T \times T$ tiles, running
all $K \times K$ taps for a tile before moving on, so the tile stays resident
in L1-D instead of being re-swept from DRAM $K \times K$ times like
\texttt{conv\_reorder}. We swept $T$ on the graded workload
(2048$\times$2048, $K=3$):

\begin{table}[H]
    \centering
    \begin{tabular}{rrrr}
        \toprule
        \textbf{$T$} & \textbf{Time (ms)} & \textbf{Speedup} & \textbf{L1-D MPKI} \\
        \midrule
        8  & 18.580 & 0.92$\times$ & 0.920 \\
        16 & 17.451 & 0.93$\times$ & 21.533 \\
        24 & 15.642 & 1.03$\times$ & 24.495 \\
        32 & 16.225 & 1.05$\times$ & 19.624 \\
        48 & 15.073 & 1.16$\times$ & 16.421 \\
        \textbf{64} & \textbf{14.809} & \textbf{1.10$\times$} & \textbf{14.148} \\
        80 & 15.479 & 1.06$\times$ & 9.852 \\
        \bottomrule
    \end{tabular}
    \caption{\texttt{conv\_tile} tile-size sweep, pinned to P-core CPU 0.}
    \label{tab:tile-sweep}
\end{table}

$T=64$ is fastest and fits the 48\,KB L1-D ($\approx 33.4$\,KB working set),
so it's what's committed in \texttt{conv\_tile.cpp}.

\begin{insightbox}
    Fastest ($T=64$) isn't the same as fewest misses ($T=8$, MPKI
    $\approx 0.92$): tiny tiles avoid conflict misses but pay for it in loop
    overhead, T from 16 to 48 all have high MPKI but $T=64$ is the best tradeoff between the two. $T=80$ has an even
    lower MPKI, but its $\approx 52$\,KB working set no longer fits the
    48\,KB L1-D, so the capacity misses it introduces cost more in real
    latency than the MPKI ratio alone shows, making it slower than $T=64$
    despite ``looking'' more cache-efficient.
\end{insightbox}

\subsection{Metrics and Plots}

We swept image size $\in \{128, 256, 512, 1024, 2048, 4096\}$ ($K=3$) for
three tile sizes -- $T=8$ (best MPKI), $T=24$ (worst, mid conflict-miss zone),
and $T=64$ (committed) -- against the naive baseline:

\begin{figure}[H]
    \centering
    \includegraphics[width=0.75\textwidth]{task1_tile_mpki_vs_size.png}
    \caption{L1-D MPKI vs.\ matrix size, one line per tile size.}
    \label{fig:tile-mpki-vs-size}
\end{figure}
\begin{figure}[H]
    \centering
    \includegraphics[width=0.75\textwidth]{task1_tile_speedup_vs_size.png}
    \caption{Speedup vs.\ matrix size, one line per tile size.}
    \label{fig:tile-speedup-vs-size}
\end{figure}

\subsection{Questions}
\begin{enumerate}[leftmargin=1.6em, itemsep=0.3em]

    \item \textbf{MPKI change naive $\to$ tiled:} The MPKI, as we increase the tile size, gets worse than the naive implementation.
    This could be because the naive implementation is already pretty cache-optimal at $K=3$. It benefits from long, contiguous horizontal memory accesses,
    which the prefetcher could easily predict and optimize. By using tiling, these contiguous accesses are broken into short burst accesses.
    Moreover, the working set for the naive approach probably already fits well into the L1D cache.

    \item \textbf{MPKI vs.\ matrix size/tile size:} For $T=8$, L1-D MPKI remains consistently low across all matrix sizes.
    For $T=24$ to $T=64$, MPKI starts low but spikes drastically as the matrix grows and finally plateaus.
    This is because, at the smallest matrix size, the entire working set fits into L1-D cache so misses are rare.
    As the matrix grows, the data exceeds L1-D capacity. The $T=8$ tile size keeps its active working set small enough to remain in cache.
    But larger tile sizes introduce vertical memory jumps that exceed cache capacity, causing misses and moreover, it isn't prefetcher friendly.

    \item \textbf{Speedup achieved:} Speedup was achieved only for larger tile sizes (> 24) with the sweet spot being $T=64$.
    The performance improved despite high L1-D misses, indicating that tiling successfully improved temporal locality was improved which led to lower latency even
    if L1-D misses made it fetch data from LLC. This could be because we don't have to go to L3 or DRAM to fetch something we use in a "short" span of time
    because we reuse that data often (temporal locality). Also the number of instructions of tiling is about half
    that of the naive implementation so that is also a factor which contributes to the speedup.
    One of the ways to make this even more efficient could be loop unrolling (to abuse spacial locality further).
\end{enumerate}

\end{document}
